% =========================================================
%  sections/04-sub-cpmk.tex
%  Bagian 4 — Sub-CPMK dan Indikator Pencapaian
% =========================================================

\section{Sub-CPMK dan Indikator Pencapaian}

Sub-CPMK merupakan penjabaran setiap CPMK menjadi kemampuan yang lebih
spesifik, terukur, dan dapat diuji. Setiap Sub-CPMK dilengkapi ranah taksonomi
Bloom, CPMK induk, dan CPL yang didukung. Indikator pencapaian dirumuskan
sebagai perilaku yang dapat diamati pada pekerjaan mahasiswa (unjuk kerja,
kode program, atau laporan).

\newcommand{\ind}[1]{\textbf{(#1)}}

\subsection{Sub-CPMK Turunan \CPMK{1} --- \istilah{Game Loop} dan Fisika Sederhana}

\begin{table}[htbp]
\centering
\caption{Sub-CPMK dan indikator pencapaian untuk \CPMK{1}}
\label{tab:subcpmk1}
\begin{tabularx}{\textwidth}{@{}L{1.5cm} L{3.9cm} Y L{1.5cm} L{2.4cm}@{}}
\toprule
\headrow \textbf{Kode} & \textbf{Rumusan Sub-CPMK} & \textbf{Indikator Pencapaian} & \textbf{Bloom} & \textbf{CPL} \\
\midrule
1.1 &
Menjelaskan struktur \istilah{game loop} (input, \istilah{update}, \istilah{render}) dan peran manajemen waktu dalam permainan. &
\ind{a} menjelaskan urutan tahapan input--\istilah{update}--\istilah{render} dengan benar;
\ind{b} mengidentifikasi peran \code{pygame.time.Clock} dan pembatasan FPS;
\ind{c} membedakan \istilah{fixed timestep} dan \istilah{variable timestep} secara konseptual. &
C2 & \CPL{03} \\
\addlinespace
1.2 &
Merancang dan mengimplementasikan gerak objek berbasis \istilah{delta time} dan vektor pada Pygame. &
\ind{a} menulis kode gerak menggunakan \code{pygame.Vector2} dan \istilah{delta time};
\ind{b} menunjukkan kecepatan objek tetap konsisten pada 30 FPS dan 60 FPS;
\ind{c} menyinkronkan posisi bertipe \istilah{float} dengan \code{Rect} bertipe bilangan bulat. &
C3 & \CPL{03} \\
\addlinespace
1.3 &
Menerapkan logika fisika sederhana (gravitasi, lompatan, pantulan) dalam simulasi 2D. &
\ind{a} menerapkan integrasi Euler untuk memperbarui kecepatan dan posisi;
\ind{b} mengimplementasikan lompatan yang hanya aktif ketika objek berada di tanah;
\ind{c} menghasilkan pantulan dengan koefisien restitusi sederhana. &
C3 & \CPL{03} \\
\addlinespace
1.4 &
Mendeteksi dan menangani \istilah{collision} menggunakan \code{Rect} dan/atau \istilah{mask} pada \istilah{sprite}. &
\ind{a} memilih metode \istilah{collision} yang tepat sesuai kebutuhan;
\ind{b} menyelesaikan tabrakan per sumbu sehingga objek tidak menembus dinding;
\ind{c} menganalisis dan memperbaiki kasus tabrakan yang salah (tumpang tindih atau lolos deteksi). &
C4 & \CPL{03}, \CPL{04} \\
\bottomrule
\end{tabularx}
\end{table}

\subsection{Sub-CPMK Turunan \CPMK{2} --- AI Dasar untuk NPC}

\begin{table}[htbp]
\centering
\caption{Sub-CPMK dan indikator pencapaian untuk \CPMK{2}}
\label{tab:subcpmk2}
\begin{tabularx}{\textwidth}{@{}L{1.5cm} L{3.9cm} Y L{1.5cm} L{2.4cm}@{}}
\toprule
\headrow \textbf{Kode} & \textbf{Rumusan Sub-CPMK} & \textbf{Indikator Pencapaian} & \textbf{Bloom} & \textbf{CPL} \\
\midrule
2.1 &
Merancang perilaku NPC sederhana (misalnya \istilah{chase}) berdasarkan posisi pemain. &
\ind{a} menghitung vektor arah NPC menuju pemain dan menormalisasinya;
\ind{b} membatasi besar kecepatan agar gerak NPC konsisten;
\ind{c} menjelaskan perbedaan AI reaktif dan AI berbasis perencanaan. &
C3 & \CPL{03} \\
\addlinespace
2.2 &
Mengimplementasikan \istilah{Finite State Machine} (idle, patrol, chase, attack/flee) untuk NPC. &
\ind{a} menggambar diagram \istilah{state}--transisi NPC;
\ind{b} mengimplementasikan transisi berbasis jarak dan \istilah{cooldown};
\ind{c} menguji setiap \istilah{state} dan transisinya secara terpisah. &
C3 & \CPL{03} \\
\addlinespace
2.3 &
Menerapkan \istilah{pathfinding} berbasis grid dengan algoritma A*. &
\ind{a} membangun representasi grid/graf berbiaya dari peta berhalangan;
\ind{b} menganalisis pengaruh heuristik terhadap optimalitas dan kinerja pencarian;
\ind{c} merekonstruksi dan memvisualisasikan jalur hasil A*. &
C4 & \CPL{03} \\
\addlinespace
2.4 &
Mengevaluasi perilaku NPC hasil kombinasi FSM dan A* terhadap skenario permainan. &
\ind{a} menyusun skenario uji dan mencatat hasil pengamatan;
\ind{b} menilai kualitas perilaku NPC (kelancaran, daya tanggap, beban komputasi);
\ind{c} merumuskan perbaikan berdasarkan hasil evaluasi. &
C5 & \CPL{03}, \CPL{04} \\
\bottomrule
\end{tabularx}
\end{table}

\subsection{Sub-CPMK Turunan \CPMK{3} --- \istilah{Game Engine} dan Pembangunan Level}

\begin{table}[htbp]
\centering
\caption{Sub-CPMK dan indikator pencapaian untuk \CPMK{3}}
\label{tab:subcpmk3}
\begin{tabularx}{\textwidth}{@{}L{1.5cm} L{3.9cm} Y L{1.5cm} L{2.4cm}@{}}
\toprule
\headrow \textbf{Kode} & \textbf{Rumusan Sub-CPMK} & \textbf{Indikator Pencapaian} & \textbf{Bloom} & \textbf{CPL} \\
\midrule
3.1 &
Menavigasi antarmuka Unity dan menyusun \istilah{project}, \istilah{scene}, serta \istilah{GameObject} dasar. &
\ind{a} mengoperasikan \istilah{Hierarchy}, \istilah{Inspector}, \istilah{Project window}, dan \istilah{Scene};
\ind{b} membuat dan menyimpan \istilah{scene} berisi kamera serta objek primitif;
\ind{c} menjelaskan padanan konsep Unity dengan konsep Pygame yang telah dipelajari. &
C3 & \CPL{04} \\
\addlinespace
3.2 &
Membangun level 2D/3D sederhana dengan \istilah{collider}, \istilah{Rigidbody}, dan kontrol pemain. &
\ind{a} menyusun level menggunakan \istilah{tilemap}/\istilah{greyboxing};
\ind{b} mengonfigurasi \istilah{Collider} dan \istilah{Rigidbody} sesuai kebutuhan mekanik;
\ind{c} menganalisis dan memperbaiki masalah fisika (\istilah{jitter}, objek menembus \istilah{collider}, jatuh tanpa henti). &
C4 & \CPL{04} \\
\bottomrule
\end{tabularx}
\end{table}

\subsection{Sub-CPMK Turunan \CPMK{4} --- Integrasi Aset Audio dan Visual}

\begin{table}[htbp]
\centering
\caption{Sub-CPMK dan indikator pencapaian untuk \CPMK{4}}
\label{tab:subcpmk4}
\begin{tabularx}{\textwidth}{@{}L{1.5cm} L{3.9cm} Y L{1.5cm} L{2.4cm}@{}}
\toprule
\headrow \textbf{Kode} & \textbf{Rumusan Sub-CPMK} & \textbf{Indikator Pencapaian} & \textbf{Bloom} & \textbf{CPL} \\
\midrule
4.1 &
Mengimpor dan mengonfigurasi aset visual (\istilah{sprite}/model, animasi, UI) ke dalam \istilah{scene}. &
\ind{a} mengimpor aset dan mengatur \istilah{importer setting} dengan tepat;
\ind{b} menyiapkan \istilah{Animator} beserta transisi animasi (misalnya idle dan run);
\ind{c} menampilkan elemen UI (skor, nyawa) yang berfungsi pada \istilah{scene}. &
C3 & \CPL{04} \\
\addlinespace
4.2 &
Mengintegrasikan aset audio dan memicu interaksi audio-visual berdasarkan peristiwa \istilah{game}. &
\ind{a} menghubungkan \code{AudioSource} dengan peristiwa \istilah{gameplay};
\ind{b} membedakan penggunaan BGM dan SFX secara tepat;
\ind{c} menyesuaikan \istilah{timing} dan level volume agar umpan balik terasa tepat. &
C4 & \CPL{04} \\
\addlinespace
4.3 &
Mengevaluasi prototipe \istilah{game} melalui \istilah{playtesting} dan memperbaiki integrasi level serta aset. &
\ind{a} menyusun \istilah{checklist} fungsional dan lembar \istilah{playtest};
\ind{b} mengidentifikasi \istilah{bug} dan masalah integrasi dari hasil \istilah{playtest};
\ind{c} menyajikan laporan perbaikan dan refleksi pemilihan teknologi. &
C5 & \CPL{04} \\
\bottomrule
\end{tabularx}
\end{table}

\subsection{Matriks Sub-CPMK terhadap CPMK}

\begin{table}[htbp]
\centering
\caption{Matriks keterkaitan Sub-CPMK dengan CPMK}
\label{tab:matriks-subcpmk}
\begin{tabularx}{\textwidth}{@{}L{3.6cm}@{\extracolsep{\fill}}C{2.6cm} C{2.6cm} C{2.6cm} C{2.6cm}@{}}
\toprule
\headrow \textbf{Sub-CPMK} & \textbf{\CPMK{1}} & \textbf{\CPMK{2}} & \textbf{\CPMK{3}} & \textbf{\CPMK{4}} \\
\midrule
1.1--1.4 & \textbf{X} & & & \\
2.1--2.4 & & \textbf{X} & & \\
3.1--3.2 & & & \textbf{X} & \\
4.1--4.3 & & & & \textbf{X} \\
\bottomrule
\end{tabularx}
\end{table}
